% Pursuing Stacks (cotes 134-2 to 134-8): archive cote 134-5, site page 22 (PDF page 23), typescript page 209.
% typescript. Draft compared: ps4.tex:998-1055.
% AI-transcribed by Claude Opus 5.5 (claude-opus-5-5, adaptive thinking, effort high), Message Batches API, batch ps2 (msgbatch_01NVwr1Yuo8qJf5rykeGBXPK), 2026-11-12; not yet checked by eye.
\scanpage{134-5/22}
\mspage{209}
the similar name of ``globular complexes''). Roughly speaking, it can be
said that the three types of complexes correspond to three ``series'' of
\uline{regular} \uline{cellular} \uline{subdivisions} of all spheres $S_n$ \deleted{\unsure{,n}} where the two-dimensio-%
nal pieces are respectively (by increasing order of ``intricacy'') bigons,
triangles, and squares. It should'nt be hard to show these are the only
series of regular cellular subdivisions of all spheres (one in each dimensi-%
on), such that for any cell of such subdivision, the induced subdivision
of the bounding sphere should still be in the series (up to isomorphism).
The existence moreover of suitable ``degeneracy'' maps, which merit a care-%
ful general definition in this context of cellular subdivisions of spheres,
is an important \inserted[typed]{common extra} \xout{striking} feature of the three basic contexts, whose exact
significance I have not quite understood still. \hand{\illegible{1}} To come back to Dold-%
Puppe, sure enough it is still valid in the hemispherical context. Wri-%
ting down the equivalence of categories in explicit terms comes out with
baffling simplicity. I wrote it down without even looking for it, in a
letter to Ronnie Brown, while explaining in a PS the ``yoga'' of associating
to a chain complex of abelian groups an $\infty$-groupoid with additive struc-%
ture (a so-called ``Picard\xout{xx} category'' but within the context of $\infty$-catego-%
ries or $\infty$-groupoids\hand{,} rather than usual categories).

On the other hand, taking multico\inserted{m}plexes instead of simple ones, \xout{whi}
which still can be interpreted as working in a category $A\hat{}$ for suitable
test category $A$ (namely, a product category of categories of the type
\hand{$\Box$}, \hand{$\triangle$} and \hand{$\Diamond$}), it is clear that Dold-Puppe's theorem \deleted{is no longer true}
as originally stated\hand{\unsure{,}} \inserted{is no longer true \deleted{\illegible{1}} these;} in such case the category of abelian group objects
of $A\hat{}$ is equivalent to a category of \uline{multiple} chain complexes. This shows
that definitely, among all possible elementary modelizers, the three in
our trilogy are distinguished indeed, as giving rise to a Dold-Puppe
theorem. A thorough understanding of this theorem would imply, I feel,
\inserted{an understanding} of which exactly are the modelizers, or elementary modelizers at any rate,
giving rise to such a theorem. I would'nt be too surprised if it turned
out that the three we got are the only ones, up to equivalence.

Another common feature of these modelizers, is that they allow for
a sweeping computational description of cohomology (or homology) invariants
in terms of the so-called\hand{``}boundary operations\hand{''}. This is visibly connected
to the \hand{(}strongly intuitive\hand{)} t\over{y}{i}e between these kind of models, \over{\illegible{1}}{and} cellular
subdivisions of spheres. However, at this level (unlike the Dold-Puppe
story) the regularity feature of the cellular subdivisions we got, and
the fact that we allow for just one (up to isomorphism) in each dimension,

\margin{22}
